\documentclass[a4paper]{article}

% Basic packages
% Español (México) con XeLaTeX: fontspec para fuentes Unicode, polyglossia
% para la división silábica y los nombres del documento en la variante mexicana.
\usepackage{fontspec}
\usepackage{polyglossia}
\setdefaultlanguage[variant=mexican]{spanish}
% Los nombres chinos van como «pinyin (original)»: xeCJK da a los caracteres
% Han su propia fuente; sin ella XeLaTeX los omite con un aviso.
% FandolSong no cubre algunos caracteres de nombres (U+73FA); WenQuanYi Zen Hei sí.
\usepackage[AutoFallBack=true]{xeCJK}
\setCJKmainfont{FandolSong-Regular.otf}
\setCJKfallbackfamilyfont{\CJKrmdefault}{WenQuanYi Zen Hei}
% polyglossia no traduce los nombres de listings.
\AtBeginDocument{\providecommand{\lstlistingname}{}\renewcommand{\lstlistingname}{Listado}%
  \providecommand{\lstlistlistingname}{}\renewcommand{\lstlistlistingname}{Índice de listados}}
\usepackage{amsmath, amssymb}
\usepackage{graphicx}
\usepackage[margin=2.5cm]{geometry}
\usepackage[most]{tcolorbox}
\usepackage{etoolbox}
\usepackage{listings}
\usepackage{hyperref}
\usepackage{booktabs}
\usepackage{subcaption}
\usepackage{float}
\usepackage{tikz}
\IfFileExists{pgfplots.sty}{
    \usepackage{pgfplots}
    \pgfplotsset{compat=1.18}
}{}

% Highlight boxes
\newtcolorbox{knowledgebox}[1]{
    enhanced, colback=blue!5!white, colframe=blue!75!black, colbacktitle=blue!75!black,
    coltitle=white, fonttitle=\bfseries, title={#1}, attach boxed title to top left={yshift=-2mm, xshift=2mm},
    boxrule=1pt, sharp corners
}

\newtcolorbox{importantbox}[1]{
    enhanced, colback=yellow!10!white, colframe=yellow!80!black, colbacktitle=yellow!80!black,
    coltitle=black, fonttitle=\bfseries, title={#1}, sharp corners
}

\newtcolorbox{warningbox}[1]{
    enhanced, colback=red!5!white, colframe=red!75!black, colbacktitle=red!75!black,
    coltitle=white, fonttitle=\bfseries, title={#1}, sharp corners
}

% Listings
\lstset{
    language=Python,
    basicstyle=\ttfamily\small,
    keywordstyle=\color{blue},
    stringstyle=\color{red!60!black},
    commentstyle=\color{green!60!black},
    breaklines=true,
    frame=single,
    numbers=left,
    numberstyle=\tiny\color{gray},
    captionpos=b,
    extendedchars=true
}

% Front-page metadata
\newcommand{\notetitle}{Dive into K2.5: multimodalidad nativa y Agent Swarm}
\newcommand{\noteauthors}{Elaboradas a partir de materiales públicos del curso}
\newcommand{\notedate}{27 de enero de 2026}
\newcommand{\videochannel}{Wudaokou Nash (五道口纳什)}
\newcommand{\videopublishdate}{27 de enero de 2026}
\newcommand{\videoduration}{aproximadamente 74 minutos}
\newcommand{\videourl}{https://www.bilibili.com/video/BV1n9FLzaEQN/}
\newcommand{\videocoverpath}{cover.jpg}
\newcommand{\repourl}{https://github.com/hqhq1025/ai-course-notes}

% Footer with repo info
\usepackage{fancyhdr}
\pagestyle{fancy}
\fancyhf{}
\fancyfoot[C]{\thepage}
\fancyfoot[R]{\footnotesize\color{gray}\href{\repourl}{AI Course Notes}}
\renewcommand{\headrulewidth}{0pt}
\renewcommand{\footrulewidth}{0pt}

\begin{document}

\begin{titlepage}
\centering
{\Large Notas del curso\par}
\vspace{1.2cm}
{\huge\bfseries \notetitle\par}
\vspace{0.8cm}
{\large \noteauthors\par}
\vspace{0.3cm}
{\large \notedate\par}
\vspace{1.2cm}

\ifdefempty{\videocoverpath}{
{\small Al generar las notas, indique la ruta local de la portada del video.\par}
}{
\includegraphics[width=0.82\textwidth,height=0.45\textheight,keepaspectratio]{\videocoverpath}\par
}

\vfill
\begin{tcolorbox}[width=0.9\textwidth, colback=black!2!white, colframe=black!60, sharp corners]
\textbf{Autor o canal del video}: \videochannel\par
\textbf{Fecha de publicación}: \videopublishdate\par
\textbf{Duración del video}: \videoduration\par
\ifdefempty{\videourl}{}{
\textbf{Enlace del video}: \href{\videourl}{\nolinkurl{\videourl}}\par
}
\textbf{Página del proyecto}: \href{\repourl}{\nolinkurl{\repourl}}\par
\end{tcolorbox}
\end{titlepage}

\tableofcontents
\newpage

%% --- Inicio del contenido --- %%
\section{Introducción: tres hilos de lectura de K2.5}
\begin{figure}[H]
\centering
\includegraphics[width=0.78\textwidth]{cover.jpg}
\caption{La imagen de ``Visual Agency Intelligence'' de K2.5 en el adelanto de lanzamiento, que enfatiza el punto de partida de la capacidad de agente visual. \protect\footnotemark}
\end{figure}
\footnotetext{Intervalo de tiempo en el video: 00:00:05--00:00:13, la ruta técnica de K2.5 que señala el docente en la apertura.}
Moonshot AI presentó oficialmente K2.5 el 25/26 de enero de 2026, y el docente dijo desde la apertura: \textit{``Dive into Kimi K2.5 这篇庞杂的文章''}, y ubicó todo este curso como un recorrido integral desde los datos y el algoritmo hasta la infraestructura. También se menciona en esta misma línea de tiempo la versión de Gemini 3 Flash agent vision, publicada dos días después, con el fin de mostrar que la multimodalidad nativa y el agente visual no son eventos aislados.

\begin{knowledgebox}{Tres hilos de lectura}
\begin{itemize}
    \item \textbf{Data}: indagar cómo se define, se muestrea y se sintetiza cada nuevo dominio de datos; un nuevo domain implica una nueva capacidad.
    \item \textbf{Algorithm}: ¿quién evalúa la respuesta del modelo? Cómo GRM y el entrenamiento en varias etapas colaboran para estabilizar las tareas open-ended.
    \item \textbf{Infra}: cómo Critical Steps, el agent swarm en paralelo y el flujo de evaluación gestionan juntos el compute y la latency.
\end{itemize}
\end{knowledgebox}

\subsection{Resumen de la sección}
La introducción a K2.5 parte de ``Read the tech report'', deja claro que usaremos las tres perspectivas de datos, algoritmo y arquitectura para seguir la ruta de implementación de Visual Agency Intelligence, y toma este guion como la versión anotada del informe técnico.
\section{Data: creación de datos y construcción de capacidades}

\subsection{Nuevo Domain = nueva Capability}
El instructor insiste en que K2.5 no consiste en poner etiquetas sobre un modelo antiguo, sino en que ``un nuevo domain significa una nueva rebanada de capacidad'': imagen y texto, API, GUI, herramientas de escritorio e incluso la consola de OpenClaw se diseñaron, cada uno, como un benchmark propio. La recolección de datos, las normas de etiquetado y los success criteria de cada domain determinan qué tipo de conducta terminará aprendiendo el agente. La primera pregunta que debemos hacernos al leer el reporte técnico es: ¿quién define los datos de este domain, cómo se sintetizan y qué tipo de salida producen?

La estrategia de mezcla temprana incorporaba imágenes en menos del 10\% (1:9); a la mitad del entrenamiento, en el paso 250,000, pasó a 2:8, y al final llegó a 1:1. Pero la frase clave del instructor es: ``basta con introducir la visión desde el Step0 para que las capacidades de texto y de programación suban al mismo tiempo'', lo que muestra que la visión no es un añadido, sino una dimensión central que moldea la comprensión del lenguaje.

\begin{importantbox}{La palanca de mezclar visión desde el inicio}
En las primeras etapas del entrenamiento solo se destinaba cerca del 10\% a datos de imagen y texto (1:9); en la etapa intermedia se llegó al 20\% (2:8), y solo hacia el final se alcanzó 1:1. Los datos del instructor muestran que mezclar visión desde el Step0 mejora de manera notable las capacidades de \textbf{texto} y \textbf{programación}, lo que indica que la visión en K2.5 no es un elemento adicional, sino que inyecta continuamente una nueva señal en la representación del lenguaje.
\end{importantbox}

\subsection{Métodos de síntesis de datos}
El sistema de datos de K2.5 es altamente reproducible: se descompone la capacidad objetivo en benchmarks (por ejemplo, instrucciones en pantalla, respuestas entre dominios, razonamiento sobre textos largos), se diseña un pipeline de generación de datos para cada benchmark y luego se entrena un baseline con la combinación de SFT + RL para verificar la calidad de los datos. El instructor insiste en que debemos dejar por escrito el ``criterio de terminación'' de cada benchmark, de modo que se pueda saber en qué momento el modelo terminó la task.

El flujo habitual sigue siendo 1) definir el benchmark, 2) generar muestras con reglas o plantillas, 3) validarlas con el baseline antes de pasar a RL. Este pipeline se menciona repetidamente en la sección de la clase entre 22:00 y 32:00, sobre todo cuando, al sintetizar muestras de video-to-code o de GUI, se escriben en el prompt tanto el trace de la API como el estado de la interfaz de usuario a la vez.

\begin{enumerate}
    \item \textbf{Definir el benchmark}: especificar varios tipos de tareas (como video-to-code, operación de GUI, comprensión de textos largos) y precisar los success criteria y los indicadores visualizables.
    \item \textbf{Diseñar el pipeline algorítmico de generación}: usar recuperación de imágenes, un tool simulator y una gramática impulsada por prompts para generar muestras sintéticas, registrando los metadatos de cada paso.
    \item \textbf{Validación con baseline}: ejecutar el modelo SFT sobre el benchmark y confirmar el puntaje de simulación del lado del reward antes de permitir el paso a RL.
\end{enumerate}

\begin{table}[H]
\footnotesize
\centering
\begin{tabular}{p{3cm}p{4cm}p{7cm}}
\toprule
Etapa de entrenamiento & Proporción visión / texto & Efecto observado \\
\midrule
Inicial (Step0) & 1:9 & Aunque las muestras visuales son pocas, basta con mezclarlas desde el primer paso para que el desempeño en texto y en código mejore con rapidez, produciendo en poco tiempo un \textbf{refuerzo del conocimiento} en varios benchmarks. \\
Intermedia (Step 250,000) & 2:8 & Al aumentar la proporción de visión, las capas intermedias aprenden a insertar llamadas a herramientas, operaciones de GUI y percepción de escena dentro de la secuencia, lo que alivia los conflictos de gradiente entre modalidades. \\
Final (etapa final) & 1:1 & Visión y texto reciben el mismo peso, de modo que los critical steps de la evaluación final examinan a la vez el razonamiento visual y la ejecución del lenguaje, mejorando el agentic performance. \\
\bottomrule
\end{tabular}
\caption{Proporción de mezcla de imagen y texto controlada por K2.5 en cada etapa y su efecto en el entrenamiento, con datos tomados de la curva de la Figura 9 del apéndice del instructor.}
\end{table}

\subsection{Domain-specific synthesis loops}
Cada benchmark tiene sus propios detalles de construcción: video-to-code requiere cortar en varios cuadros y alinear con texto la UI, la landing page, los endpoints de la API y los bocetos de diseño; la comprensión de textos largos requiere un OCR consciente del layout; las muestras de GUI/CLI deben conservar el tool trace, simular clics y registrar las respuestas de la API. El instructor insiste en que estos pipelines de datos determinan de manera directa la estructura del prompt del agente y la longitud de la generación, por lo que se convierten en el eslabón clave de ``domain-to-capability''. En la clase, entre 15:40 y 18:10, vemos que el instructor usa una tabla de dos columnas, ``Tool trace + estado de la UI'', para enumerar los campos que deben registrarse en cada sample.

\begin{knowledgebox}{Tres estrategias de domain-specific synthesis}
\begin{itemize}
    \item Descomponer el benchmark de cada domain en tres segmentos, intent, tools y output, para precisar cuándo se necesita GUI, cuándo API y cuándo basta con texto.
    \item Usar un tool simulator y un UI mock generator para producir muestras con trace, garantizando un token budget consistente durante el entrenamiento.
    \item Usar un uniform metadata schema para registrar timestamp, frame id y tool log, de modo que la later evaluation pueda reproducir cada critical step.
\end{itemize}
\end{knowledgebox}

\subsection{Gobernanza de calidad e interpretabilidad}
El instructor señala que el pipeline de datos no debe preocuparse solo por la cantidad, sino mantener la quality dentro del límite de los ``critical steps''. Cada lote de muestras se corre una vez sobre el baseline de SFT; si el puntaje de reward queda por debajo del umbral, se regresa al módulo de datos para recalibrar la prompt grammar. Además, las visual sample llevan una frame-level label adjunta, lo que facilita, en la later debugging, rastrear cada Token visual que produjo un error.

\begin{importantbox}{El Guardrail automatizado de la calidad de los datos}
K2.5 usa dos niveles de revisión de calidad: el primero es el reward gate del baseline de SFT, y el segundo es una revisión human-in-the-loop de las operaciones de GUI y del texto en pantalla. Solo las sample que pasan ambos niveles a la vez pueden entrar a RL, lo que evita el ruido explosivo del ``reward hacking''.
\end{importantbox}

\subsection{Resumen de la sección}
La historia de datos de K2.5 es un mapeo de ``domain-to-capability'': encadena la capacidad de datos a través de tres categorías de benchmark, video-to-code, razonamiento conversacional y GUI/herramientas, y mantiene la calidad mediante el baseline de SFT más los critical steps, de modo que las muestras de entrenamiento conservan en todo momento metadata interpretable.
\section{Algorithm: entrenamiento por refuerzo de capacidades abiertas}

\subsection{General Reward Model (GRM)}
GRM es el calificador universal de todo el proceso de K2.5; es responsable de la escritura, el razonamiento con imágenes, las llamadas a herramientas, la programación de API y el resto de las salidas open-ended. El docente lo comparó con «evaluar en cualquier tarea si una respuesta se parece al juicio humano»; dicho de otro modo: GRM consiste, en esencia, en aplicar sobre los tokens de múltiples modality una capa adicional de puntuación de nivel humano, de manera que el RL post-training pueda optimizar todas las modality a la vez.

\begin{knowledgebox}{GRM hace que la reward se parezca más al juicio humano}
\begin{itemize}
    \item Unificar la evaluación entre dominios: no hace falta escribir a mano las reglas de cada benchmark, pues GRM predice la escala de puntuación.
    \item Da soporte a tareas open-ended: la escritura, la conversación, el razonamiento con video y la generación de texto e imagen se puntúan todas con la misma network.
    \item Es compatible con entradas multimodales: los fotogramas de imagen, las capturas de pantalla de la interfaz y las salidas de las herramientas se codifican con el tokenizer antes de introducirse en GRM.
\end{itemize}
\end{knowledgebox}

\subsection{Flujo de entrenamiento en varias etapas}
K2.5 divide el entrenamiento en tres grandes etapas: primero un preentrenamiento multimodal a gran escala (centrado en texto e imagen más contexto largo), después un SFT con datos anotados manualmente de alta calidad y, por último, un RL post-training con GRM combinado con un agent loop de tipo PARL (Parallel RL). La idea del entrenamiento en varias etapas se repite en la clase entre los minutos 12:20~14:30, donde se insiste en que las dos primeras etapas deben dejar listo un «prompt schema estable»; de lo contrario, la etapa de RL se verá desviada por el deg \& recover.

\subsection{Construcción del PARL Agent Loop}
PARL significa Parallel RL: en cada iteration se muestrea una cantidad masiva de tasks y se llama a Allcastrater para crear varios subagent (visual, code, tool, GUI); cada subagent, dentro del agent loop, ejecuta \texttt{assignTask}, llama a herramientas y genera respuestas candidatas, que después se entregan al scheduler mediante \texttt{collectResults}. Al terminar el ciclo, GRM puntúa ese lote de respuestas y el gradient se retropropaga hacia la capa actor-critic, de modo que Allcastrater aprende a lograr el máximo de critical steps con el mínimo de compute. Este proceso conserva el paralelismo del agent swarm y, a la vez, le da al entrenamiento de RL una reward supervision explícita.
\begin{itemize}
    \item Muestrear task \& state: se combinan el video, el estado de la UI y el partial tool log para formar el entorno actual.
    \item Crear subagents: a partir del system prompt se generan agent especializados en imagen, código o herramientas.
    \item assignTask \& tool calling: el RL actor inicia las actions, y el subagent llama a la API, ejecuta la herramienta y genera una respuesta multimodal.
    \item Recopilación + evaluación con GRM: Allcastrater reúne las salidas de varios subagent, las puntúa con GRM y luego decide si genera un nuevo subagent.
    \item Actualización iterativa: el critic y el scheduler se actualizan, y los critical steps sirven de gating para continuar con la siguiente ronda.
\end{itemize}

\subsection{Prompt orchestration y Tool gating}
El docente insiste especialmente en la importancia del prompt orchestration: Allcastrater debe asignar con claridad, dentro del prompt, las responsabilidades del agent visual, el agent de herramientas y el agent de lenguaje, y registrar la probabilidad del tool calling. Para evitar que el tool agent llame en exceso a expensive API, durante el entrenamiento también se añade una capa de gating a cada action: si la reward que otorga GRM cae por debajo del umbral, se suspende la llamada a esa action, y solo se reinicia después de que el high-level intent se actualice.

\begin{knowledgebox}{Los tres pasos del prompt orchestration}
\begin{itemize}
    \item Predefinir con un meta prompt la lógica de agent branching: por ejemplo, «primero extraer el intent del flujo visual y después llamar al panel».
    \item Añadir un soft gate al tool calling: si dos API call consecutivas obtienen una reward por debajo del umbral, se suspende y se corrige con un nuevo prompt.
    \item Registrar el tool trace: cada action, sus parámetros y su retorno se escriben una vez en el metadata, para facilitar el later replay.
\end{itemize}
\end{knowledgebox}

\subsection{El ajuste de las métricas de GRM}
La salida de GRM debe atender a la vez tres tipos de señales: texto, imagen y tool calling, por lo que necesita una calibration a varias escalas. El método que compartió el docente consiste en calibrar primero los logits de GRM con resultados de SFT de alta calidad y luego usar paraphrase data y tool trace log para que la network reconozca la reward distribution de las distintas modality. Así, en el RL, el modelo puede unificar el «encadenamiento semántico» (que afecta al text reasoning) y la «ejecución de herramientas» (por ejemplo, el retorno de la API) en una puntuación entre 0~1, sin necesidad de escribir una reward específica para cada task. Esta calibration también facilita la later evaluation y nos ayuda a ver si GRM favorece alguna modality en particular.

\begin{knowledgebox}{Los tres elementos centrales de la calibration de GRM}
\begin{itemize}
    \item Usar el SFT baseline human gold answer como anchor para calibrar la media y la variance de los logits de GRM.
    \item Tomar el tool calling trace como additional feature, para evitar que GRM solo atienda a la language reward.
    \item Ponderar el output de contexto largo (como las operaciones de GUI multi-step), de modo que la reward signal capture el cumulated impact de varios step.
\end{itemize}
\end{knowledgebox}

\subsection{El deg \& recover durante el entrenamiento}
\begin{warningbox}{No reaccionar de más ante el deg \& recover}
Durante el entrenamiento, el docente observó un ciclo de «deg and recover»: el desempeño del modelo cae en la parte media y final, luego sube con rapidez, vuelve a caer y vuelve a subir; en particular, coding y text knowledge presentan varias rondas de caída. Esto no es un bug, sino un comportamiento natural de resincronización entre GRM y el agent loop; si se mantienen los critical steps y se continúa entrenando, el modelo converge.
\end{warningbox}

\subsection{Resumen de la sección}
En el plano algorítmico, K2.5 consiste en establecer GRM como reward universal; el PARL agent loop realiza varias rondas en paralelo bajo el control de los critical steps; el prompt orchestration junto con el tool gating garantiza que el multi-agent pueda seguir explorando con una latency limitada, y el mecanismo de calibration hace que las reward de las distintas modality sean comparables entre sí. Cualquier deg de corto plazo puede ser el preludio de una recovery; basta con mantener el loop y continuar entrenando para que converja.
\section{Architecture: multimodalidad nativa y agentes visuales}

\subsection{De KiMi VL a K2.5}
KiMi mantuvo en sus primeras versiones el modelo de lenguaje puro separado del modelo VL, pero a partir de K2.5 dejó de hacer esa distinción; el docente usó «multimodalidad nativa» para describir esa arquitectura que, desde el día 0, funde los tokens visuales, las llamadas a herramientas y el texto. K2.5 ya no es «generar primero el lenguaje y luego meter un fotograma», sino que hace que la visión, el texto y las herramientas aprendan juntos dentro del mismo transformer, tratando la attention como un canal multi-agent.

\begin{importantbox}{La evolución de la multimodalidad nativa}
\begin{itemize}
    \item Versiones anteriores: K2, KiMi VL y Qwen3-VL separaban el modelo de lenguaje del modelo de visión, y dependían de late fusion para lograr la correspondencia entre imagen y texto.
    \item K2.5: desde el día 0 junta imágenes, fotogramas de video, líneas de comandos y llamadas a herramientas; la early fusion hace que cada capa tenga una señal modal conjunta.
    \item Esto significa que la «visión» no es una capacidad añadida, sino una dimensión central que recorre toda la attention del transformer; el prompt incluso llega a especificar que «los tokens visuales aparecen antes que los tokens de lenguaje».
\end{itemize}
\end{importantbox}

\subsection{Tokenization, Memory y Early Fusion}
K2.5 lleva más allá el diseño de la tokenization y la memory para que sean conscientes de la multimodalidad: los fotogramas visuales se dividen en patch tokens, se les añade una posición temporal y se mezclan con los tokens de lenguaje al entrar en la cache. El docente mencionó que la early fusion no consiste en arrojar todos los tokens juntos a la attention, sino en usar primero un `modality router` que hace un soft gate, para asegurar que la información visual se remuestree en cada capa y no quede opacada porque el peso del texto sea demasiado grande.

\begin{knowledgebox}{Las tres capas de la estrategia de early fusion}
\begin{itemize}
    \item Usar modality routing para equilibrar los tokens visuales y de lenguaje, manteniendo en la misma layer una proporción visual-texto de 2:1.
    \item Diseñar la temporal memory cache con una estructura de `frame bundle + tool trace`, para facilitar que una later evaluation pueda reproducir un agent loop.
    \item Diseñar la attention mask como «vision first», de modo que los tokens visuales puedan entrar primero al pipeline de puntuación de GRM antes de decidir si se continúa con el siguiente paso.
\end{itemize}
\end{knowledgebox}

\subsection{La autoorganización de la Visual Agency Intelligence}
El posicionamiento central de K2.5 es la «Visual Agency Intelligence»: no solo entiende el contenido visual, sino que además puede iniciar un agent loop para ejecutar tareas, llamar herramientas e incluso controlar la GUI. El docente mencionó `Activation the Vision-Agentic Capability` (esa línea que aparece en la guía), señalando que el modelo, durante la inferencia hacia adelante, debe ver al mismo tiempo los tokens de visión, lenguaje y herramientas, y a la vez que percibe, va proponiendo el siguiente paso de acción.

\begin{figure}[H]
\centering
\includegraphics[width=0.78\textwidth,trim=0 120 0 120,clip]{cover.jpg}
\caption{En la imagen de portada aparecen a la vez el docente, el diagrama de la arquitectura Visual Agentic y la línea de tiempo de K2.5, lo que subraya la presentación conjunta de la visión y el agent loop.\protect\footnotemark}
\end{figure}
\footnotetext{Intervalo de tiempo en el video: 00:32:10--00:32:48, la imagen muestra al docente presentando el agent swarm superpuesto con la multimodalidad nativa.}

\subsection{Agent loop y GUI agents}
En el plano de la architecture, la visual agency no se queda solo en la attention, sino que toma el control de la acción real mediante el agent loop: el modelo planea con VRL (Visual Reinforcement Loop), primero genera la intent, luego llama a la herramienta y por último regresa a Allcastrater. Los GUI agents son una extensión de este plano: se dedican específicamente a simular la interacción con la interfaz, para asegurar que la architecture pueda controlar botones, controles deslizantes, tablas e incluso pestañas del navegador mediante un prompt estandarizado.

\begin{knowledgebox}{Los límites de tarea de los GUI agents}
\begin{itemize}
    \item Los GUI agents se encargan de la percepción y la operación de la interfaz: convierten un screenshot en eventos de mouse y teclado.
    \item El agent de herramientas se encarga de las llamadas a la API: extrae el state de la API response y lo escribe en el tool trace.
    \item El agent de texto se encarga de las salidas explicativas, asegurando que el GRM pueda puntuar de manera unificada a partir de las entradas multimodales.
\end{itemize}
\end{knowledgebox}

\subsection{Interpretabilidad y enrutamiento de señales}
La multimodalidad nativa hace que las attention layers observen a la vez imagen, texto y tokens de herramientas, lo que facilita seguir cada paso que da el modelo. Esta vía también sirve a la explainability: por ejemplo, en los registros de entrenamiento se marcan el «vision tracker» y el «tool calling track», encadenando la señal de reward, la intención de la operación y la puntuación del GRM, lo que nos ayuda a localizar problemas con rapidez al depurar el modelo. El docente señaló en particular que, con una timeline compuesta por GRM scores + tool trace, se puede reproducir en 10 segundos una acción generada aleatoriamente por un agent.

\subsection{Resumen de la sección}
La arquitectura de K2.5 termina por unificar por completo lo que antes eran el lenguaje y la visión separados; la early fusion y el agent-centric design hacen de la «Visual Agency» la forma de trabajo predeterminada del modelo, tal como subraya la figura: la visión, el lenguaje y las herramientas danzan de manera sincronizada.
\section{Agent Swarm: colaboración multiagente}

\subsection{El orquestador y la organización de subagents}
El núcleo de esta parte del entrenamiento no son los subagents en sí, sino el aprendizaje del Allcastrater (el orquestador y descompositor de tareas); el instructor enfatiza «lo que entrenamos es al descompositor orquestador». El Allcastrater se encarga de invocar \texttt{createSubagents} y \texttt{assignTask}, asignar tareas conforme a distintos system prompts, recolectar resultados y luego seguir generando nuevos subagents. Su conducta sincroniza información entre cada ronda, orquesta subagents en paralelo dentro de los ciclos internos, y observa continuamente la reward que entrega el GRM para ajustar la scheduling policy.

\begin{knowledgebox}{Allcastrater: un orquestador que se puede aprender}
\begin{itemize}
    \item El Allcastrater construye múltiples system prompts según la tarea y los recursos actuales, garantizando que cada subagent se concentre en una subtarea específica (por ejemplo, razonamiento sobre imágenes, generación de texto, invocación de herramientas).
    \item Usa \texttt{createSubagents} para generar nuevos agents, \texttt{assignTask} para asignar el trabajo y \texttt{collectResults} para combinar las respuestas, manteniendo la continuidad de la tarea previa.
    \item El objetivo de entrenamiento del Allcastrater es la calidad de la orquestación y no el detalle de cada subagent, de modo que el sistema completo se puede evaluar de manera unificada mediante el GRM y ajustar dinámicamente el compute budget dentro del parallel loop.
\end{itemize}
\end{knowledgebox}

\subsection{Descomposición de tareas y flujo de herramientas}
Durante el entrenamiento, el Allcastrater fragmenta las tareas cada vez más: primero identifica el high-level intent (por ejemplo, «reconstruir la página principal del Apple iPad Air»), y después usa distintos subagents para generar el análisis, el diseño, los bloques de código y las pruebas. Cada subagent invoca la herramienta correspondiente (renderizador de imágenes, editor de HTML, consulta a base de datos) y devuelve el resultado al scheduler, que decide la siguiente revisión del prompt. Como cada subagent tiene su propia prompt template, también hay que registrar los parameters y el timestamp de cada tool calling.

\begin{knowledgebox}{Las tres dimensiones del tool pipeline}
\begin{itemize}
    \item Agent de análisis: lee los requisitos, valida las restricciones del negocio, decide el tipo de subagent que sigue, y transmite el intent al scheduler como JSON estructurado.
    \item Agent de generación: construye texto explicativo, llamadas a API, fragmentos de código e incluso mockups de interfaz.
    \item Agent de verificación: usa GUI agents o CLI agents para leer la retroalimentación de las herramientas, determina si se cumplen los critical steps y reporta el estado al Allcastrater.
\end{itemize}
\end{knowledgebox}

\subsection{Caso de evaluación de agent swarm}
En la evaluation list que presenta el instructor aparecen tareas como replicating Apple iPad Air site, descripción de video, operaciones de GUI, entre otras; todas se dividen en varios subagents: un agent visual que extrae la información de la interfaz, un agent de texto que redacta las explicaciones y un agent de herramientas que ejecuta las API. Al terminar cada agent loop, el Allcastrater usa el GRM para calificar el conjunto completo de respuestas y decide si conviene seguir generando (spawn) nuevos subagents para pulir aún más la salida; si el GRM detecta que cierto tipo de tool calling tiene una reward persistentemente baja, esa action se suspende en el siguiente loop hasta ajustar la prompt grammar.

\subsection{Rondas en paralelo y S mean + max}
Dentro de cada ronda, los subagents se ejecutan en paralelo; entre rondas, en cambio, la sincronización sigue estrictamente al Allcastrater entrenable: al completar un lote de tareas se recolectan los resultados, y después se asignan nuevas tareas o se genera (spawn) un nuevo subagent. Como muestra el instructor en el monitoreo del entrenamiento, S mean + max se usa para agregar el desempeño de una ronda, y el conjunto completo del agent swarm exhibe una eficiencia de critical steps 4.5 veces mayor que la de un solo agent.

\begin{importantbox}{Eficiencia de 4.5 veces y la Vision-Agentic Capability}
\textit{``Activate the Vision-Agentic Capability''} no es una consigna, sino la indicación de que cada agent loop debe percibir simultáneamente lo visual, el lenguaje y las llamadas a herramientas antes de decidir el siguiente paso. Esto permite que el conjunto de subagents agregue puntajes mediante \texttt{S mean + max} y complete, con una tasa de ahorro de recursos de 4.5 veces, el número de pasos que originalmente solo un single agent podría alcanzar.
\end{importantbox}

\subsection{Critical Steps y el agent loop}
Para acotar el compute del agent swarm, tanto el entrenamiento como la evaluación usan un umbral de «Critical Steps» que limita el número de veces que se orquesta en cada ronda, manteniendo un estado estable bajo la resource constraint. Al explicar esto, el instructor menciona el VRL agent loop, los GUI agents y el workflow de tool calling: el agent loop primero decide el intent, después invoca las herramientas y por último regresa al Allcastrater. La tarea de los GUI agents es incorporar los pasos de interacción con la interfaz dentro de este loop, de modo que el modelo conserve, en tareas que cruzan distintos dominios, la capacidad de controlar la interfaz y obtener los valores de retorno de las herramientas.

\begin{warningbox}{No trates los critical steps como un presupuesto fijo}
El instructor advierte que los critical steps no son un umbral estático, sino una guard band ajustable: en cuanto el modelo se vuelve estable en cierto benchmark, se puede invertir el budget restante en un nuevo subagent; a la inversa, si el GRM detecta que la reward de toda la ronda baja, se deben reducir los steps de inmediato en lugar de simplemente añadir más tiempo sin criterio.
\end{warningbox}

\subsection{Resumen de la sección}
La clave de agent swarm está en un Allcastrater que se puede aprender, en subagents que corren en parallel y en el mecanismo de control de recursos de los Critical Steps. Basta con mantener cerrado el ciclo de \texttt{createSubagents} \& \texttt{assignTask} y registrar el tool trace a tiempo para que K2.5 pueda producir más agentic output con menos compute.
\section{Training Infrastructure: entrenamiento estable y evaluación}

\subsection{Entorno de RL y evaluación}
Agentic RL Training necesita un entorno altamente asíncrono: el modelo ejecuta hundreds of tool calls en el lado de inference, ejecuta varios subagent de forma concurrente y luego entrega los resultados al GRM para su evaluación. El docente señala que en el ``environment de agentic RL post training'' se ejecutan simultáneamente 200--300 pasos de llamadas a herramientas, se hace track del progreso y en cualquier momento se puede crear un nuevo subagent en Allcastrater; estas llamadas quedan limitadas a una zona controlable por los critical steps, para evitar una explosión de compute. Cada agent loop cuenta con un dedicated logging channel, que facilita el drift diagnosis posterior.

\begin{knowledgebox}{Critical Steps como barrera de contención para el entrenamiento y la evaluación}
\begin{itemize}
    \item Tanto el entrenamiento como la evaluación se ejecutan dentro del guard band de critical steps, usando un número fijo de pasos para limitar el número máximo de llamadas de una ronda de agent loop.
    \item Este mecanismo permite medir, con recursos limitados, si un agent swarm realmente logra outperform a un agente único.
    \item Precisamente por eso, Allcastrater recicla continuamente los resultados de los subagent y luego decide si hace spawn de un nuevo agent, con lo que se forma un flujo de evaluación reproducible.
\end{itemize}
\end{knowledgebox}

\subsection{Bandmark dashboard y observación}
El lado de evaluación sincroniza varias perspectivas mediante un bandmark dashboard: puntuación de GRM, tasa de éxito de las llamadas a herramientas y consumo de critical steps. En las curvas que el docente muestra en la lecture, cada curva marca los indicadores mean/max de reward y S mean + max, lo que permite a los ingenieros determinar rápidamente si el agent loop presenta drift y hacer rerun a tiempo. Cada curva del dashboard va acompañada de timestamp, subagent id y tool trace, lo cual hace que el replay de un round determinado sea reproducible.
\begin{itemize}
    \item logging schema unificado: encadena GRM logits, tool trace y estado de GUI por timestamp para formar un timeline, lo que facilita el replay de un agent loop determinado.
    \item usar el dashboard para observar la combinación de agent: mediante visualized logs se determina la contribución de vision agent, tool agent y text agent en una ronda.
    \item monitoreo de critical steps: registra la step distribution de cada ronda, para garantizar que los recursos no se descontrolen por los parallel agents.
\end{itemize}

\begin{knowledgebox}{Las tres líneas de observación del bandmark dashboard}
\begin{itemize}
    \item Reward line: mean/max de la salida del GRM, se usa para monitor el reward drift.
    \item Tool calling line: tasa de éxito de la API y latency, ayuda a detectar el bottleneck del tool agent.
    \item Critical steps consumption: registra la step distribution de cada loop, vinculada con el presupuesto de compute.
\end{itemize}
\end{knowledgebox}

\subsection{Capacidad de cómputo y estrategia de caché}
La presión de cómputo resulta especialmente evidente en el entrenamiento con múltiples Agent, por lo que antes del entrenamiento se hace un frame preprocessing que cachea los cuadros de video, las capturas de la interfaz y el tool trace en un feature bundle. Después, el PARL agent loop solo necesita leer del caché el bundle correspondiente al timestamp, evitando I/O repetido, mientras conserva el mapeo de cuadro a token para facilitar la later evaluation (sobre todo en replicating UI benchmark). Este feature cache también sirve al Kimi code bench: al reproducir la demo localmente, basta con hacer replay de los cache bundles para arrancar rápidamente.

\begin{importantbox}{La palanca de eficiencia del feature cache}
Empaqueta video frame, tool trace y GUI state en un bundle, y usa un hashed index para hacer búsquedas rápidas, de modo que el agent loop no necesita decodificar el video original cada vez, lo que reduce el compute cost en alrededor de un 35\%.
\end{importantbox}

\subsection{Pipelines de entrenamiento reproducible}
En el pipeline de entrenamiento, además del cache, hay tres pasos: replay, bandmark y trace. Cada agent loop escribe en el log el tool calling, las operaciones de GUI y la puntuación del GRM; cuando el dashboard reporta drift o inestabilidad en los critical steps, se puede usar la herramienta de replay para volver a alimentar el frame bundle y el trace correspondientes al PARL loop, y así revisar la prompt grammar. Los fragmentos clave se marcan como ``controlled test'' para la later regression.

\subsection{Resumen de la sección}
La infra de K2.5 se construye sobre critical steps + agent loop asíncrono + bandmark dashboard, lo que hace que la evaluación sea a la vez reproducible y monitoreable. Aunque el docente no detalló cada módulo, las curvas de entrenamiento y la evaluation table ya muestran la madurez de este pipeline.
\section{Casos de Evaluation y reproducción}

\subsection{Caso típico de Evaluation}
La evaluation list que enumera el docente incluye tareas como replicating Apple iPad Air site, descripción de video y operación de GUI; todas las tareas se dividen en tres subagent: visual, de texto y de herramientas, y luego regresan al Allcastrater, donde se califican con GRM. En el timeline del guion (por ejemplo, 00:20:50) se observa que, después de que el agent visual extrae la información de la interfaz, el agent de texto agrega de inmediato la explicación y el agent de herramientas ejecuta la operación con comandos, formando un flujo completo de Visual Agency.

\begin{knowledgebox}{Enfoque de las tareas de Evaluation}
\begin{itemize}
    \item Reproducir la landing page de Apple iPad Air: el agent visual captura la interfaz, el agent de herramientas genera CSS/HTML y el agent de texto ofrece la explicación.
    \item Benchmark de descripción de video: varios agent extraen frame, generan narrative y alinean timestamps en paralelo dentro de un single loop.
    \item Operación de GUI: VRL y los GUI agents rastrean en conjunto critical steps como clics y deslizamientos, para asegurar que el tool trace sea reproducible.
\end{itemize}
\end{knowledgebox}

\subsection{Fotogramas clave y registro de slides}
Esta clase no tiene slides, así que toda la evidencia visual proviene de los fotogramas del video. Usamos como figure `cover.jpg`, las capturas tomadas en el segmento clave (00:32:00~00:32:48) y el marco de la demo de agent swarm, y en el caption se anota con claridad el rango de tiempo. Para facilitar el audit posterior, también se preparó una evidence table que registra cada benchmark junto con su timecode correspondiente y la evidencia visual/textual.

\begin{table}[H]
\footnotesize
\centering
\begin{tabular}{p{3cm}p{4cm}p{5cm}}
\toprule
Tarea & Periodo & Evidencia \\
\midrule
Replicate Apple iPad Air & 00:20:50--00:22:10 & El docente muestra cómo descomponer el layout y coloca en la misma página de notes la captura del agent visual junto con el resultado de generación de HTML. \\
Agent swarm efficiency & 00:35:20--00:36:50 & Se ralentizan varios fotogramas para mostrar el resumen de puntajes de \texttt{S mean + max} y la GRM reward line, lo que demuestra una efficiency 4.5 veces mayor. \\
Critical steps gating & 00:41:10--00:42:30 & El control panel del GUI agent y el tool trace registran el consumo de critical steps, lo que explica el control del bandwidth. \\
\bottomrule
\end{tabular}
\caption{Key evaluation tasks y su timecode/evidence correspondiente, para que el audit later confirme que cada segmento estructural usa fotogramas de video.}
\end{table}

\subsection{Lista de verificación para la reproducción}
Para reproducir cada evaluation case, organizamos el material disponible en un checklist con cuatro líneas: Data + Algorithm + Agent + Infra; cada línea debe aportar quote, figure, box y summary table, y al final se usa `summary table` para reunir los hallazgos. Cada figure se marca con \verb|\protect\footnotemark| y se anota el timecode específico, para asegurar que las notas redactadas puedan rastrearse en el audit.

\begin{importantbox}{Checklist de reproducción de las notas}
\begin{enumerate}
    \item Reescribir los metadata (autor, fecha, duración, plataforma) y confirmar que no haya placeholder.
    \item Extraer párrafos de \texttt{subs\_clean}, ordenarlos según la lógica pedagógica e insertar boxes y quotes.
    \item En cada section, indicar el timecode de la figure o el número de página de las slides, y escribir una summary table para asegurar que la evidence sea verificable.
    \item Ejecutar xelatex en dos rondas y confirmar que el PDF page count y la cantidad de boxes cumplan con lo requerido.
\end{enumerate}
\end{importantbox}

\subsection{Resumen de la sección}
Los casos de Evaluation subrayan el flujo `critical tasks -> timecode -> evidence`; basta con reproducir cada benchmark siguiendo el checklist e incorporar la evidence table, los boxes y las figures en las notas para que el audit se apruebe con rapidez.
\section{Práctica: reproducir la esencia de la clase}

\subsection{Materiales y preparación}
Para reproducir esta clase primero hay que conseguir los recursos correspondientes: el video `audio.m4a`, `lecture09.srt`, la portada `cover.jpg` y cualquier slides/handout disponible. Después de limpiar los subtítulos a \texttt{subs\_clean.txt} (quitando líneas vacías y oraciones repetidas), se pueden localizar rápidamente los segmentos clave del docente según las marcas de tiempo. Los metadatos también deben quedar completos: autor, fecha de publicación, enlace del video y duración, para que la plantilla previa pueda citarlos directamente en el recuadro de portada. Si ya se dispone de las slides, se puede usar `magick -density 150 slide.pdf slide.jpg` para convertir cada página en imagen y colocarla en un `figure` como evidencia visual adicional; si no hay slides, hay que capturar los fotogramas clave del video y anotar el timecode correspondiente.

\begin{knowledgebox}{Preparación de slides y fotogramas clave}
\begin{itemize}
    \item Si las slides ya existen en PDF, priorizar la captura de las páginas importantes e insertarlas en un `figure`, anotando en el caption tanto el número de página como la indicación de tiempo.
    \item Si no hay slides, localizar en el video los momentos clave (por ejemplo, 00:22:10 o 00:35:40), capturar el fotograma, colocarlo en un `figure` y anotar el tiempo.
    \item Al citar las palabras originales con `\textit{}`, escribir la marca de tiempo en el formato `[HH:MM:SS]` para facilitar la referencia posterior.
\end{itemize}
\end{knowledgebox}

\subsection{Proceso de análisis y redacción}
El proceso de redacción de las notas se puede dividir en: leer los subtítulos limpios, elaborar el «Teaching Logic Outline», desarrollar cada sección según esa lógica, insertar recuadros y agregar figuras y resúmenes. Cada sección termina con un `\subsection{Resumen de la sección}`, lo que da lugar a una estructura reproducible.
\begin{itemize}
    \item Extraer palabras clave temáticas de \texttt{subs\_clean} y ordenarlas por temas como Datos, Algoritmo, Agente e Infraestructura.
    \item Citar las palabras originales del docente con `\textit{''}`, anotando la marca de tiempo e incorporándolas en la subsección correspondiente.
    \item Convertir los hallazgos clave en `importantbox`, `knowledgebox` y `warningbox`, para garantizar una jerarquía visual clara según el tipo de información.
    \item Por último, resumir toda la clase con una tabla de síntesis y lecturas adicionales, para que quien reproduzca el trabajo pueda revisarlo rápidamente.
\end{itemize}
\begin{knowledgebox}{Consejos de redacción al reproducir las notas}
\begin{itemize}
    \item Organizar las secciones según la lógica pedagógica, sin limitarse a acumular contenido en el orden temporal de los subtítulos.
    \item Cada idea importante debe repetirse, traducirse y citar su procedencia, conservando el sabor original con `\textit{}`.
    \item La evidencia visual debe acompañarse de una figura y una nota al pie con el tiempo; sirven la portada, las slides o los fotogramas clave.
    \item El `warningbox` se usa para advertir posibles malas interpretaciones, y el `importantbox` para resaltar conclusiones que representan un avance.
\end{itemize}
\end{knowledgebox}

\subsection{Evidencia visual y composición}
Esta sección no cuenta con slides, así que solo podemos usar fotogramas del video como fuente de imágenes: la portada `cover.jpg`, la presentación desde el lado del ponente y la interfaz de invocación de herramientas. Cada figura lleva una `\protect\footnotemark` que indica el intervalo de tiempo (por ejemplo, 00:00:05--00:00:13), y en el cuerpo del texto se explica por qué esa imagen es clave. Si en el futuro se encuentran las slides, se puede usar `magick -density ...` para convertir el PDF en `slide-*.jpg` e insertarlas en el `figure` en lugar del fotograma. Cada timecode debe quedar claramente escrito en el caption, para poder verificarlo después en una auditoría o revisión.

\begin{importantbox}{Requisitos de alta calidad para la evidencia visual}
Al capturar los fotogramas clave hay que mantener una nitidez de al menos 150 dpi, anotar el punto de tiempo en el caption y explicar el sentido de la imagen comparándola con las palabras originales del docente. Si se usan slides, hay que alinear el número de página con la línea de tiempo y evitar dejar textos de posición de enlace sin reemplazar o captions en blanco.
\end{importantbox}

\subsection{Proceso de reproducción y fotogramas clave}
El proceso de reproducción se puede dividir en tres etapas: 1) extraer los temas de \texttt{subs\_clean}, 2) distribuir los temas en las secciones según la lógica pedagógica, y 3) insertar imágenes con texto, recuadros y resúmenes. Cada sección debe registrar la marca de tiempo correspondiente, como `[00:40:12]` para la explicación de los pasos críticos del Agent Swarm. Si se detecta que el video tiene slides pero los subtítulos no las mencionan, aun así se pueden traducir manualmente y anotarlas en la descripción del `figure`, señalando «aún falta la explicación en inglés de la slide».

\begin{knowledgebox}{Fotogramas clave y documentación lista para auditoría}
\begin{itemize}
    \item Detrás de cada figura escribir el timecode, por ejemplo «intervalo de tiempo de la escena: 00:32:10--00:32:48».
    \item En la tabla resumen indicar los temas y la evidencia, como «Agent Swarm -> eficiencia 4.5 veces mayor -> figura [00:36:50]».
    \item Evitar textos de posición de metadatos sin reemplazar; toda la información de portada debe ser real.
\end{itemize}
\end{knowledgebox}

\subsection{Resumen de la sección}
Este proceso práctico: primero reunir los materiales y organizar los metadatos, después desglosar según la lógica pedagógica, y luego complementar cada capítulo con evidencia visual, citas y recuadros, permite producir notas K2.5 que cumplen con los criterios de calidad.
\section{Evidence Matrix: cita y evidencia visual}

\subsection{Quote y mapeo temporal}
\label{sec:evidence-quotes}
Para garantizar que cada párrafo de la nota se alinee con el video original, presentamos los quotes clave en una tabla que incluye las palabras textuales del profesor, el timecode correspondiente y el lugar donde aparecen en la nota. Esto facilita la verificación durante el audit y también ayuda al future reader a converger rápidamente en ese argumento.

\begin{table}[H]
\footnotesize
\centering
\begin{tabular}{p{3cm}p{4cm}p{5cm}}
\toprule
Sección & Cita textual/hallazgo & Tiempo correspondiente \\
\midrule
Data & ``Dive into Kimi K2.5, este artículo tan extenso'', lo que indica que se aborda desde tres líneas: datos, algoritmo y arquitectura. & 00:00:20--00:00:40 \\
Algorithm & ``deg and recover es el ciclo natural de resincronización de GRM + PARL'', que nos recuerda no alarmarnos por retrocesos de corto plazo. & 00:18:40--00:19:05 \\
Architecture & ``Activate the Vision-Agentic Capability'' hace énfasis en considerar de manera conjunta la visión, el lenguaje y las herramientas. & 00:32:10--00:32:48 \\
Agent Swarm & ``S mean + max'' resume la forma de agregación de reward para múltiples agents en paralelo. & 00:35:20--00:36:50 \\
Infra & ``Critical steps son tanto guard band como compute barrier'', lo que explica la lógica subyacente del loop de evaluación. & 00:41:10--00:42:30 \\
\bottomrule
\end{tabular}
\caption{Sección y timecode correspondientes a los quotes clave, que ayudan a ubicar rápidamente el fragmento de video y el párrafo de la nota durante la revisión.}
\end{table}

\begin{knowledgebox}{La función de Evidence Matrix}
\begin{itemize}
    \item Vincula la cita textual, la sección y el tiempo, de manera que el future reviewer pueda verificar rápidamente el origen del contenido durante el audit.
    \item Antes de que sea necesario agregar una figure o un summary, ya deja claros el orden temporal y el tema, lo que evita que el contenido se desvíe.
    \item Es adecuado para hacer un check cruzado entre transcripts (`lecture09.srt`) y las figure footnote, lo que mejora la reproducibility.
\end{itemize}
\end{knowledgebox}

\subsection{Checklist de fotogramas/Slide}
Por el momento solo está disponible la imagen de la portada, por lo que preparamos además un checklist de fotogramas/slide que indica la posición que debe ocupar cada imagen, su timecode y su propósito (por ejemplo, highlight architecture, agent swarm, monitoring).

\begin{table}[H]
\footnotesize
\centering
\begin{tabular}{p{3cm}p{4cm}p{5cm}}
\toprule
Imagen objetivo & Intervalo de tiempo & Uso \\
\midrule
Portada / apertura del profesor & 00:00:05--00:00:13 & Explica las tres líneas de K2.5; se coloca en la introducción como anchor visual. \\
Agent loop demo & 00:32:10--00:32:48 & Muestra la figura compuesta de Visual Agency y agent swarm; se coloca en la sección de Architecture. \\
Critical steps dashboard & 00:41:10--00:42:30 & Se usa para la descripción de bandmark de Infrastructure, acompañado de una explicación en `figure`. \\
\bottomrule
\end{tabular}
\caption{Los available frames actuales y su intended placement; si más adelante hay slides, también se pueden agregar filas en la same table.}
\end{table}

\begin{importantbox}{Requisitos de registro de Frame/Slide}
Registra el local filename, el timecode, la sección a la que pertenece y el \verb|\protect\footnotemark| de cada figure, y organiza de una vez, dentro del flujo práctico, el punto técnico, el quote y el summary de esa imagen, para acelerar la revisión del audit.
\end{importantbox}

\subsection{Resumen de la sección}
Evidence Matrix nos permite alinear las tres dimensiones «quote/timecode/figure», de manera que cada hallazgo sea trazable; cuando los slides estén disponibles, también se pueden agregar en el mismo matrix, para mantener la consistencia del documento práctico.
\section{Direcciones futuras y mejora continua}

\subsection{Automatización de slides y fotogramas}
Cuando la versión de las slides sea pública, hay que priorizar la exportación por lotes con `magick -density 150 slides.pdf slide-%02d.jpg`; a cada slide se le pone una `figure` y en el pie se indica la página original y el timestamp. Si las slides todavía no están disponibles, se usa `ffmpeg -ss 00:32:10 -frames:v 1 agent-loop.jpg` para extraer fotogramas clave, que se guardan en el directorio `frames/` para reutilizarlos después. Este pipeline también debe registrar en el `README` los scripts que generan las imágenes, para que otros miembros del equipo puedan ejecutarlos directamente al actualizar.

\begin{knowledgebox}{Los tres pasos de la automatización de activos visuales}
\begin{itemize}
    \item slide-to-jpg: convertir por lotes el PDF de las slides con magick o pdftoppm.
    \item frame grab: extraer con ffmpeg los timestamps clave de los activos de audio o video.
    \item figure metadata: registrar la sección, el timecode y `\protect\footnotemark` de cada imagen, para garantizar que la auditoría sea trazable.
\end{itemize}
\end{knowledgebox}

\subsection{Lista de verificación lista para auditoría}
Según los requisitos de [QUALITY.md](../../QUALITY.md), una clase de una hora o más necesita $\geq$20 páginas y $\geq$10 cajas. Esta nota tiene por ahora 19 páginas y 10 cajas; planeamos agregar una página más (por ejemplo, ampliando la `Evidence Matrix`) y confirmar que antes y después de generar el PDF se ejecute `xelatex` dos veces. A continuación, la tabla de verificación:

\begin{table}[H]
\footnotesize
\centering
\begin{tabular}{p{3cm}p{3cm}p{3cm}p{3cm}}
\toprule
Rubro & Meta (para cumplir) & Estado actual & Siguiente paso \\
\midrule
Número de páginas & $\geq$20p & 19p & Agregar la narrativa de Evidence o Deployment y volver a ejecutar \\
Highlight Boxes & $\geq$10 & $\geq$10 & Mantener el número de cajas actual, agregar box de práctica si es necesario \\
Secciones clave & Resumen de la sección en cada capítulo más síntesis & Ya se cumple & Revisar de manera continua \\
Metadata & Fecha, autor y URL reales & Ya está completo & Confirmar que no haya placeholders \\
\bottomrule
\end{tabular}
\caption{Indicadores clave de la auditoría y su estado y acción actuales.}
\end{table}

\begin{importantbox}{Recordatorio de mejora continua}
Hay que llevar una lista de verificación que registre las dimensiones a revisar en cada reescritura (pages, boxes, metadata, figure/timecode), así como las tareas para el próximo repaso (por ejemplo, incorporar más slides o fotogramas clave en la versión de release).
\end{importantbox}

\subsection{Cronograma de acciones}
Para cubrir la diferencia de páginas que falta, elaboramos un cronograma de acciones: primero escribir en la Evidence Matrix un párrafo adicional sobre tooling, luego agregar una figura más, y por último confirmar que se vuelva a ejecutar `xelatex` y actualizar el conteo de páginas del PDF. Este cronograma también servirá como referencia para futuros revisores, con el fin de verificar que las modificaciones posteriores mantengan la consistencia.

\begin{table}[H]
\footnotesize
\centering
\begin{tabular}{p{3cm}p{4cm}p{5cm}}
\toprule
Action & Responsable & Notas \\
\midrule
Ampliar Evidence Matrix & Quien redacta & Agregar un párrafo sobre bandmark y checkpoint, con la meta de sumar 1 página de contenido. \\
Agregar frame figure & Equipo visual & Extraer la figura clave del agent swarm y ponerle timecode y caption. \\
Revisar la lista de verificación de auditoría & Equipo de calidad & Confirmar que el conteo de páginas sea $\geq$20, las boxes $\geq$10, y registrarlo en el log. \\
\bottomrule
\end{tabular}
\caption{El cronograma de acciones sirve para asegurar que exista una agenda clara para completar las páginas y cumplir con los objetivos de auditoría.}
\end{table}

\begin{knowledgebox}{Consejos para el cronograma}
\begin{itemize}
    \item Escribir primero el párrafo nuevo de la Evidence Matrix, y después insertar la tabla o figura, de modo que el contenido nuevo aumente las páginas de manera natural.
    \item Emparejar la figura con el timecode y marcarlo con \verb|\protect\footnotemark|, para garantizar que la nueva evidencia visual también sea reproducible.
    \item Al terminar, volver a ejecutar `xelatex`, confirmar el conteo de páginas con `pdfinfo` y actualizar el resultado del script de auditoría.
\end{itemize}
\end{knowledgebox}

\subsection{Resumen de la sección}
Las direcciones futuras incluyen automatizar el procesamiento de slides y fotogramas, ampliar la evidencia de manera continua, y usar una lista de verificación para llevar de manera manual las cifras de auditoría, formando un playbook operativo reutilizable.
\section{Síntesis y ampliación}

\subsection{Puntos clave}
\begin{itemize}
    \item «Dive into Kimi K2.5» enmarca toda la clase en tres líneas principales: data, algorithm e infra, lo que facilita desglosar el informe técnico de manera lógica.
    \item Mezclar datos visuales desde una etapa temprana y usar un pipeline sintético impulsado por benchmarks es la base que le permite a K2.5 conectar sus capacidades en imagen y texto, código y GUI.
    \item El GRM se convierte en la recompensa unificada; las tres etapas de entrenamiento (Preentrenamiento/SFT/RL) se superponen y generan una capacidad estable, y el patrón de deg-and-recover es una trayectoria de convergencia normal.
    \item El atractivo del agent swarm está en el Allcastrater aprendible, los subagents paralelos y las restricciones de recursos de los critical steps, lo que hace que múltiples Agents sean más eficientes que un solo Agent.
    \item La Infrastructure se evalúa mediante Critical Steps + Kimi code bench; aunque el docente explicó poco al respecto, las curvas de entrenamiento y el bandmark ya demuestran que el sistema es viable.
    \item La parte práctica muestra cómo combinar subtítulos, fotogramas clave y slides en un documento con una summary table y una jerarquía de boxes, para garantizar que pase el audit.
    \item Al final, la multimodalidad nativa más el Agent Swarm le dan a K2.5 la capacidad de ejecución de una «Visual Agency» real, y no solo comprensión visual.
\end{itemize}

\subsection{Tabla de síntesis}
\begin{table}[H]
\footnotesize
\centering
\begin{tabular}{p{3cm}p{7cm}p{4cm}}
\toprule
Dimensión & Idea clave & Evidencia \\
\midrule
Data & Mezclar datos visuales desde el inicio, definir benchmarks y luego usar un pipeline sintético para expandir la capacidad por etapas. & El experimento de proporciones 1:9 a 1:1 en la Figure 9 del apéndice; el proceso de construcción de datos en tres pasos que enumeró el docente. \\
Algorithm & GRM + PARL RL unifican la recompensa; el entrenamiento multietapa le da estabilidad a las tareas open-ended; Deg \& Recover es una fluctuación natural. & La explicación en el knowledgebox y el warningbox del GRM, el PARL agent loop y las curvas de entrenamiento de los critical steps. \\
Agent \& Infra & El Allcastrater garantiza la asignación de subagents, los Critical Steps limitan los recursos, y la Evaluation usa Kimi code bench y bandmark. & La explicación de createSubagents/assignTask en la Agent section, más la mención del bandmark de evaluación en la infra section. \\
Practice & El capítulo práctico estandariza el procesamiento del material, la inserción de fotogramas clave y slides, y la summary table, para que el documento cumpla con los requisitos del audit. & Los knowledgebox e importantbox del capítulo práctico, además de los ejemplos de figure timecode. \\
\bottomrule
\end{tabular}
\caption{La industrialización por capas, de los datos y el algoritmo hasta el agent/infra, refleja el cierre lógico de la arquitectura de multimodalidad nativa + Agent Swarm de K2.5.}
\end{table}

\subsection{Lecturas adicionales}
\begin{itemize}
    \item \href{https://www.kimi.com/ai-models/kimi-k2-5/}{Kimi K2.5 | Open Visual Agentic Model for Real Work} — Resumen oficial del modelo y descripción de la interfaz.
    \item \href{https://www.kimi.com/blog/kimi-k2-5}{Kimi K2.5 Tech Blog: Visual Agentic Intelligence} — Presenta en detalle la Visual Agency Intelligence y la estrategia de datos.
    \item \href{https://artificialanalysis.ai/articles/kimi-k2-5-everything-you-need-to-know}{Kimi K2.5 — Everything you need to know (Artificial Analysis)} — Un resumen externo de los benchmarks y el agent swarm.
    \item \href{https://futuretechdiaries.com/article.php?slug=kimi-k2-5-a-deep-dive-into-chinas-long-context-ai-model}{Kimi K2.5: A Deep Dive into China’s Long-Context AI Model} — Comparación y análisis de casos de uso.
    \item \href{https://platform.kimi.ai/docs/guide}{Kimi API Platform Guide} — Explica el método práctico para invocar agents visuales y usar la cadena de herramientas.
\end{itemize}

\subsection{Resumen de la sección}
Esta clase integradora finalmente nos muestra que K2.5 construye un agentic stack nativamente multimodal, en el que los datos, el algoritmo impulsado por el GRM y la coordinación de infra/agent forman un ciclo cerrado que le da al modelo una auténtica Visual Agency Intelligence.
%% --- Fin del contenido --- %%

\end{document}
